\par\addvspace{12pt}\Needspace{9\baselineskip}

\section{线性变换*}

\label{sol:sec:vector:4}

% ITEM: vector-translation-nonlinearity
\solutionitem{vector-translation-nonlinearity}{思考：平移与线性 *}
平移\(T \left( a \right) = a + b\)（\(b \ne \vec{0}\)）为什么不是线性变换？
\solutionanswer
线性变换必须把零向量映为零向量，而这里\(T(0)=b\ne0\)，所以不是线性变换。也可直接比较：
\[
T(u+v)=u+v+b,\qquad T(u)+T(v)=u+v+2b.
\]
平移保持两点之间的差，却不保持以固定原点为基准的向量加法。


% ITEM: vector-basis-images
\solutionitem{vector-basis-images}{例3.10 *}
设\(T\)是\(\mathbb{R}^{2}\)上的线性变换，已知\(T \left( 1 , 0 \right) = \left( 1 , 2 \right)\)，\(T \left( 0 , 1 \right) = \left( 3 , - 1 \right)\)，求\(T \left( 2 , 1 \right)\)与\(T\)在标准基下的矩阵。
\solutionanswer
由线性性，
\[
T(2,1)=2T(1,0)+T(0,1)=(5,3).
\]
把两个标准基向量的像按列排列，矩阵为
\[
A=\begin{bmatrix}1&3\\2&-1\end{bmatrix}.
\]
验证时，\(A(2,1)^{\mathsf T}=(5,3)^{\mathsf T}\)。


% ITEM: vector-isometries-shear-area
\solutionitem{vector-isometries-shear-area}{思考：长度、夹角与面积 *}
考虑本节的伸缩、旋转、正交投影、镜像和错切变换。 上述变换中，哪些保持向量的长度与夹角？错切变换保持图形的面积吗？
\solutionanswer
旋转和镜像都保持长度与夹角。伸缩\((x,y)\mapsto(ax,by)\)在\(a=b\)时保持夹角，只有\(a=b=1\)时还保持全部长度。投影一般改变长度与夹角。

错切\(H_k(x,y)=(x+ky,y)\)在\(k=0\)时为恒等变换；\(k\ne0\)时，原来垂直的两个单位向量的像为\((1,0),(k,1)\)，数量积为\(k\)，后一向量的长度为\(\sqrt{1+k^2}\)，所以不保持长度与夹角。

错切保持面积。对以\(u=(x_1,y_1),v=(x_2,y_2)\)为邻边的平行四边形，其面积为\(|x_1y_2-x_2y_1|\)，而
\[
(x_1+ky_1)y_2-(x_2+ky_2)y_1=x_1y_2-x_2y_1.
\]
因此平行四边形与三角形的面积不变，再把多边形分成三角形即可。几何上，错切将每条水平截线平移而保持其长度，也说明了面积为什么不变。


% ITEM: vector-noncommuting-transforms
\solutionitem{vector-noncommuting-transforms}{例3.11 *}
将横坐标伸长为原来的\(2\)倍，纵坐标不变，再绕原点逆时针旋转90°。如果交换这两个变换的次序，结果相同吗？
\solutionanswer
记伸缩、旋转的矩阵分别为
\[
D=\begin{bmatrix}2&0\\0&1\end{bmatrix},\qquad
R=\begin{bmatrix}0&-1\\1&0\end{bmatrix}.
\]
先伸缩后旋转对应\(RD\)，交换次序对应\(DR\)：
\[
RD=\begin{bmatrix}0&-1\\2&0\end{bmatrix},\qquad
DR=\begin{bmatrix}0&-2\\1&0\end{bmatrix}.
\]
两者不同。例如\((1,0)\)分别变为\((0,2)\)和\((0,1)\)。这也核对了正文例后给出的计算。


% ITEM: vector-rotation-addition
\solutionitem{vector-rotation-addition}{思考：旋转复合与和角公式 *}
用变换的几何意义说明\(R_{\alpha} R_{\beta} = R_{\alpha + \beta}\)，并通过矩阵乘法重新得到正弦、余弦的和角公式。
\solutionanswer
先绕原点旋转\(\beta\)，再旋转\(\alpha\)，每个非零向量的方向总共改变\(\alpha+\beta\)，长度不变。因此
\[
R_\alpha R_\beta=R_{\alpha+\beta}.
\]
另一方面，记\(c_\alpha=\cos\alpha,s_\alpha=\sin\alpha\)，矩阵乘积为
\[
\begin{bmatrix}c_\alpha&-s_\alpha\\s_\alpha&c_\alpha\end{bmatrix}
\begin{bmatrix}c_\beta&-s_\beta\\s_\beta&c_\beta\end{bmatrix}
=
\begin{bmatrix}
c_\alpha c_\beta-s_\alpha s_\beta&
-(c_\alpha s_\beta+s_\alpha c_\beta)\\
s_\alpha c_\beta+c_\alpha s_\beta&
c_\alpha c_\beta-s_\alpha s_\beta
\end{bmatrix}.
\]
与\(R_{\alpha+\beta}\)的第一列比较，便得
\[
\begin{aligned}
\cos(\alpha+\beta)&=\cos\alpha\cos\beta-\sin\alpha\sin\beta,\\
\sin(\alpha+\beta)&=\sin\alpha\cos\beta+\cos\alpha\sin\beta.
\end{aligned}
\]


% ITEM: vector-projection-fibers
\solutionitem{vector-projection-fibers}{思考：投影的原像 *}
对投影到\(x\)轴的变换，给定一个像点后，原来的点有哪些？用方程组的解与核分别描述。
\solutionanswer
投影为\(T(x,y)=(x,0)\)。给定像点\((a,0)\)，原像满足方程\(x=a\)，而\(y\)任意，故全部原像为
\[
\{(a,t):t\in\mathbb R\}=(a,0)+\operatorname{span}(0,1).
\]
这里\((a,0)\)是一个特解，\(\ker T=\operatorname{span}(0,1)\)是对应齐次方程的解集。丢失的是垂直于\(x\)轴的分量。


% ITEM: vector-inverse-order
\solutionitem{vector-inverse-order}{思考：复合的逆变换 *}
若\(A , B\)都可逆，证明\(\left( B A \right)^{- 1} = A^{- 1} B^{- 1}\)，并从连续做两个变换、再恢复原向量的过程说明次序为什么如此。
\solutionanswer
由矩阵乘法的结合律，
\[
(A^{-1}B^{-1})(BA)=I,\qquad
(BA)(A^{-1}B^{-1})=I.
\]
所以\((BA)^{-1}=A^{-1}B^{-1}\)。

若先做\(A\)对应的变换，再做\(B\)对应的变换，要恢复原向量，须先撤销后做的\(B\)，再撤销\(A\)，故恢复过程的矩阵正是\(A^{-1}B^{-1}\)。


% ITEM: vector-linear-map-tests
\solutionitem{vector-linear-map-tests}{习题3.4（1） *}
判断下列映射是否为\(\mathbb{R}^{2}\)上的线性变换：\\
\exercisesubitem{1}\(T \left( x , y \right) = \left( x + 1 , y \right)\)；\\
\exercisesubitem{2}\(T \left( x , y \right) = \left( x^{2} , y \right)\)；\\
\exercisesubitem{3}\(T \left( x , y \right) = \left( 2 x - y , x + 3 y \right)\)。
\solutionanswer
\exercisesubitem{1}不是线性变换，因为\(T(0,0)=(1,0)\ne(0,0)\)。

\exercisesubitem{2}不是线性变换。例如
\[
T(2(1,0))=(4,0)\ne2T(1,0)=(2,0).
\]

\exercisesubitem{3}是线性变换，因为两个输出坐标都是输入坐标的线性组合。直接代入可得\(T(ku+sv)=kT(u)+sT(v)\)，其矩阵为
\[
\begin{bmatrix}2&-1\\1&3\end{bmatrix}.
\]


% ITEM: vector-geometric-matrices
\solutionitem{vector-geometric-matrices}{习题3.4（2） *}
写出下列变换在标准基下的矩阵，并判断是否可逆：\\
\exercisesubitem{1}横坐标伸长为原来的\(3\)倍，纵坐标不变；\\
\exercisesubitem{2}绕原点逆时针旋转90°；\\
\exercisesubitem{3}正交投影到直线\(y = x\)上；\\
\exercisesubitem{4}关于直线\(y = - x\)作镜像；\\
\exercisesubitem{5}\(x^{\prime} = x + 2 y\)，\(y^{\prime} = y\)。
\solutionanswer
各变换的矩阵依次为
\[
\begin{array}{c|c|c}
&\text{矩阵}&\text{可逆性}\\ \hline
\mathrm{(i)}&\begin{bmatrix}3&0\\0&1\end{bmatrix}&\text{可逆}\\[6pt]
\mathrm{(ii)}&\begin{bmatrix}0&-1\\1&0\end{bmatrix}&\text{可逆}\\[6pt]
\mathrm{(iii)}&\dfrac12\begin{bmatrix}1&1\\1&1\end{bmatrix}&\text{不可逆}\\[6pt]
\mathrm{(iv)}&\begin{bmatrix}0&-1\\-1&0\end{bmatrix}&\text{可逆}\\[6pt]
\mathrm{(v)}&\begin{bmatrix}1&2\\0&1\end{bmatrix}&\text{可逆}
\end{array}
\]
投影将\((1,-1)\)映为零，故不可逆。其余四个变换分别用反向伸缩、旋转\(-90^\circ\)、同一次镜像以及错切参数\(-2\)恢复。


% ITEM: vector-invertible-map-data
\solutionitem{vector-invertible-map-data}{习题3.4（3） *}
设\(e_{1} = \left( 1 , 0 \right)\)，\(e_{2} = \left( 0 , 1 \right)\)，线性变换\(T\)满足\(T \left( e_{1} \right) = \left( 2 , - 1 \right)\)，\(T \left( e_{2} \right) = \left( 1 , 1 \right)\)。求\(T \left( 3 , 2 \right)\)、\(T\)的矩阵与逆矩阵，并求满足\(T \left( x , y \right) = \left( 1 , 2 \right)\)的向量。
\solutionanswer
基向量的像给出矩阵
\[
A=\begin{bmatrix}2&1\\-1&1\end{bmatrix}.
\]
因此
\[
T(3,2)=3(2,-1)+2(1,1)=(8,-1).
\]
由于\(\det A=3\)，逆矩阵为
\[
A^{-1}=\frac13\begin{bmatrix}1&-1\\1&2\end{bmatrix}.
\]
所求向量为
\[
\begin{bmatrix}x\\y\end{bmatrix}
=A^{-1}\begin{bmatrix}1\\2\end{bmatrix}
=\begin{bmatrix}-1/3\\5/3\end{bmatrix}.
\]


% ITEM: vector-rectangular-image-kernel
\solutionitem{vector-rectangular-image-kernel}{习题3.4（4） *}
对线性映射\(T : \mathbb{R}^{3} \longrightarrow \mathbb{R}^{2}\)，\(T \left( x , y , z \right) = \left( x + z , y + z \right)\)，求矩阵、像与核，并写出方程\(T \left( x , y , z \right) = \left( 1 , 2 \right)\)的全部解。
\solutionanswer
矩阵、像与核分别为
\[
A=\begin{bmatrix}1&0&1\\0&1&1\end{bmatrix},\qquad
\operatorname{Im}T=\mathbb R^2,\qquad
\ker T=\operatorname{span}(-1,-1,1).
\]
像为整个平面，是因为\(T(a,b,0)=(a,b)\)；核的表达式来自\(x+z=y+z=0\)。

方程\(T(x,y,z)=(1,2)\)的全部解为
\[
(x,y,z)=(1,2,0)+t(-1,-1,1),\qquad t\in\mathbb R.
\]


% ITEM: vector-rotation-stretch-order
\solutionitem{vector-rotation-stretch-order}{习题3.4（5） *}
先将平面绕原点逆时针旋转\(\theta\)角，再沿两个坐标方向分别伸缩\(a , b\)倍（\(a , b > 0\)），写出复合变换的矩阵。如果交换次序，什么条件下矩阵相同？
\par
\solutionanswer
令\(D=\operatorname{diag}(a,b)\)。先旋转后伸缩的矩阵为
\[
DR_\theta=
\begin{bmatrix}a\cos\theta&-a\sin\theta\\b\sin\theta&b\cos\theta\end{bmatrix},
\]
交换次序得到
\[
R_\theta D=
\begin{bmatrix}a\cos\theta&-b\sin\theta\\a\sin\theta&b\cos\theta\end{bmatrix}.
\]
两者相同当且仅当\((a-b)\sin\theta=0\)，即\(a=b\)，或\(\theta\)为\(\pi\)的整数倍。


% ITEM: vector-projection-reflection-identities
\solutionitem{vector-projection-reflection-identities}{习题3.4（6） *}
证明：投影矩阵满足\(P_{\theta}^{2} = P_{\theta}\)，镜像矩阵满足\(M_{\theta}^{2} = I_{2}\)。从几何上解释这两个等式，并求它们各自的像与核。
\solutionanswer
取沿投影直线的单位向量\(u=(\cos\theta,\sin\theta)\)。由\(P(v)=(v\cdot u)u\)，
\[
P^2(v)=((v\cdot u)u\cdot u)u=(v\cdot u)u=P(v).
\]
故\(P_\theta^2=P_\theta\)。几何上，已经落在直线上的点再投影一次保持不变。

镜像满足\(M=2P-I\)，于是
\[
M^2=4P^2-4P+I=I.
\]
这说明连续镜像两次恢复原向量。又有
\[
\begin{aligned}
\operatorname{Im}P&=\operatorname{span}(\cos\theta,\sin\theta),\\
\ker P&=\operatorname{span}(-\sin\theta,\cos\theta),\\
\operatorname{Im}M&=\mathbb R^2,\qquad \ker M=\{0\}.
\end{aligned}
\]
镜像的像与核也可直接由\(M^{-1}=M\)得到。


% ITEM: vector-rotation-new-basis
\solutionitem{vector-rotation-new-basis}{习题3.4（7） *}
设\(f_{1} = \left( 1 , 1 \right)\)，\(f_{2} = \left( 1 , 2 \right)\)。求旋转90°的变换在这组基下的矩阵，与标准基下的矩阵比较。
\solutionanswer
旋转默认取逆时针方向。标准基下的矩阵是
\(R=\begin{bmatrix}0&-1\\1&0\end{bmatrix}\)。求新基下的矩阵，应把基向量的像也写成新基的坐标：
\[
R f_1=(-1,1)=-3f_1+2f_2,\qquad
R f_2=(-2,1)=-5f_1+3f_2.
\]
故所求矩阵为
\[
[R]_{(f_1,f_2)}=\begin{bmatrix}-3&-5\\2&3\end{bmatrix}.
\]
它与标准基下的矩阵不同，但描述同一次旋转；例如它的平方仍为\(-I\)，对应旋转\(180^\circ\)。


% ITEM: vector-nonstandard-basis-data
\solutionitem{vector-nonstandard-basis-data}{习题3.4（8） *}
设\(T\)是\(\mathbb{R}^{2}\)上的线性变换，已知\(T \left( 1 , 1 \right) = \left( 3 , 1 \right)\)，\(T \left( 1 , - 1 \right) = \left( 1 , 3 \right)\)。求\(T \left( 2 , 3 \right)\)与\(T\)在标准基下的矩阵，并判断是否可逆。
\solutionanswer
利用
\[
(1,0)=\frac{(1,1)+(1,-1)}2,\qquad
(0,1)=\frac{(1,1)-(1,-1)}2,
\]
得到\(T(1,0)=(2,2),T(0,1)=(1,-1)\)。所以
\[
A=\begin{bmatrix}2&1\\2&-1\end{bmatrix},\qquad
T(2,3)=2(2,2)+3(1,-1)=(7,1).
\]
由于\(\det A=-4\ne0\)，变换可逆。


% ITEM: vector-rank-and-solutions
\solutionitem{vector-rank-and-solutions}{习题3.4（9） *}
设\(T : \mathbb{R}^{4} \longrightarrow \mathbb{R}^{3}\)是线性映射，核的维数为\(2\)。求\(T\)的秩，判断它能否将\(\mathbb{R}^{4}\)映满\(\mathbb{R}^{3}\)。若\(T \left( X \right) = B\)有解，写出全部解的一般形式，并说明独立参数的个数。
\solutionanswer
由秩与核的维数关系，\(\operatorname{rank}T=4-2=2\)。像只有二维，因此不能映满三维空间\(\mathbb R^3\)。

若方程有解，任取一个特解\(X_0\)，再取核的一组基\(u_1,u_2\)，全部解为
\[
X=X_0+s u_1+t u_2,\qquad s,t\in\mathbb R.
\]
有两个独立参数。


% ITEM: vector-parameter-map
\solutionitem{vector-parameter-map}{习题3.4（10） *}
对线性变换
\begin{equation*}
    T \left( x , y , z \right) = \left( x + z , y + z , x + y + t z \right) ,
\end{equation*}
讨论\(t\)取何值时\(T\)可逆；不可逆时，求像与核的一组基。对不同的\(t\)，求方程\(T \left( x , y , z \right) = \left( 1 , 1 , 2 \right)\)的全部解。
\solutionanswer
记输出为\((u,v,w)\)。由前两个坐标得\(x=u-z,y=v-z\)，代入第三个坐标：
\[
w=u+v+(t-2)z.
\]
当\(t\ne2\)时，每个目标向量都有唯一原像：
\[
z=\frac{w-u-v}{t-2},\qquad x=u-z,\quad y=v-z.
\]
所以变换可逆。对目标\((1,1,2)\)，唯一解为\((1,1,0)\)。

当\(t=2\)时，输出恰好满足\(w=u+v\)，故
\[
\begin{aligned}
\operatorname{Im}T&=\operatorname{span}\bigl((1,0,1),(0,1,1)\bigr),\\
\ker T&=\operatorname{span}(-1,-1,1).
\end{aligned}
\]
列出的向量分别线性无关，因而确为各自的一组基。此时所给方程的全部解为
\[
(x,y,z)=(1,1,0)+s(-1,-1,1),\qquad s\in\mathbb R.
\]


% ITEM: vector-rotated-ellipse
\solutionitem{vector-rotated-ellipse}{习题3.4（11） *}
将单位圆\(x^{2} + y^{2} = 1\)先沿\(x\)轴伸长为原来的\(2\)倍，再绕原点逆时针旋转45°。记像点为\(\left( x^{\prime} , y^{\prime} \right)\)，求变换的矩阵与像曲线的方程，并求原点到像曲线上点的距离的最大值与最小值。
\solutionanswer
变换矩阵为
\[
A=R_{\pi/4}\begin{bmatrix}2&0\\0&1\end{bmatrix}
=\frac1{\sqrt2}\begin{bmatrix}2&-1\\2&1\end{bmatrix}.
\]
由\(x'=(2x-y)/\sqrt2,y'=(2x+y)/\sqrt2\)，反解得
\[
x=\frac{x'+y'}{2\sqrt2},\qquad y=\frac{y'-x'}{\sqrt2}.
\]
代入单位圆方程，像曲线为
\[
\frac{(x'+y')^2}{8}+\frac{(y'-x')^2}{2}=1,
\]
也就是\(5x'^2-6x'y'+5y'^2=8\)。

旋转保持到原点的距离，而伸缩后的距离平方为
\[
4x^2+y^2=1+3x^2,\qquad x^2+y^2=1.
\]
所以距离的最大值为\(2\)，最小值为\(1\)，分别在原圆的横轴端点、纵轴端点取到。


% ITEM: vector-idempotent-decomposition
\solutionitem{vector-idempotent-decomposition}{习题3.4（12） *}
设\(T\)是线性空间\(V\)上的线性变换，满足\(T^{2} = T\)。证明：对任意\(v \in V\)，\(T \left( v \right)\)属于像，而\(v - T \left( v \right)\)属于核；像与核的交集只含零向量。由此证明，每个向量都可以唯一地表示为像中的一个向量与核中的一个向量之和。
\solutionanswer
按定义，\(T(v)\in\operatorname{Im}T\)。又由\(T^2=T\)，
\[
T(v-T(v))=T(v)-T^2(v)=0,
\]
所以\(v-T(v)\in\ker T\)。因此
\[
v=T(v)+(v-T(v))
\]
给出了所需分解。

若\(w\in\operatorname{Im}T\cap\ker T\)，写成\(w=T(u)\)，则
\[
w=T(u)=T^2(u)=T(w)=0.
\]
所以交集只含零向量。若\(v=a+b=a'+b'\)是两种分解，便有
\(a-a'=b'-b\)属于该交集，从而\(a=a',b=b'\)，分解唯一。


% ITEM: vector-shear-images
\solutionitem{vector-shear-images}{习题3.4（13） *}
设\(T \left( x , y \right) = \left( x + y , y \right)\)。分别求直线\(y = x\)的像与原像的方程，并求单位圆的像方程。判断像曲线是否关于坐标轴对称。
\solutionanswer
像与原像应分别追踪。记\(T(x,y)=(X,Y)=(x+y,y)\)。

原直线上的点为\((t,t)\)，像为\((2t,t)\)，故直线的像是\(Y=X/2\)。原像则要求输出点在直线\(Y=X\)上，即\(y=x+y\)，故原像是直线\(x=0\)。

反解\(x=X-Y,y=Y\)，单位圆的像方程为
\[
(X-Y)^2+Y^2=1.
\]
它不关于任何一条坐标轴对称。例如\((1,1)\)在曲线上，而关于两坐标轴的对称点\((1,-1),(-1,1)\)均不在曲线上。


% ITEM: vector-invertible-composition
\solutionitem{vector-invertible-composition}{习题3.4（14） *}
设\(S , T\)是有限维空间\(V\)上的线性变换。证明：若\(S \circ T\)可逆，则\(S\)与\(T\)都可逆。指出证明中使用有限维条件的步骤。
\solutionanswer
若\(T(v)=0\)，则\((S\circ T)(v)=0\)。复合可逆意味着其核只含零向量，故\(T\)单射。

任给\(w\in V\)，由于复合满射，存在\(v\in V\)使\(S(T(v))=w\)，所以\(S\)满射。

这两步尚未使用有限维。有限维条件用于下一步：同一个有限维空间上的线性变换，单射与满射等价。因此\(T,S\)都可逆。


% ITEM: vector-polynomial-shift
\solutionitem{vector-polynomial-shift}{习题3.4（15） *}
在\(P_{2}\)上规定\(T \left( p \right) \left( x \right) = p \left( x + 1 \right)\)。求\(T\)在基\(1 , x , x^{2}\)下的矩阵与逆变换，并与平面上的平移比较。
\solutionanswer
对多项式的线性组合，代入\(x+1\)与相加、数乘相容，所以\(T\)线性。基向量的像为
\[
T(1)=1,\qquad T(x)=1+x,\qquad T(x^2)=1+2x+x^2.
\]
故矩阵为
\[
A=\begin{bmatrix}1&1&1\\0&1&2\\0&0&1\end{bmatrix}.
\]
逆变换是\(T^{-1}(p)(x)=p(x-1)\)，其矩阵为
\[
A^{-1}=\begin{bmatrix}1&-1&1\\0&1&-2\\0&0&1\end{bmatrix}.
\]
平面上的平移将每个位置向量加上一个固定非零向量；这里则把整个多项式作为向量，替换自变量不会给零多项式增加常数，也保持多项式的线性组合。二者所作用的对象不同。


% ITEM: vector-trigonometric-shift
\solutionitem{vector-trigonometric-shift}{习题3.4（16） *}
在函数空间\(V = \operatorname{span}\left( \sin x , \cos x \right)\)上规定\(T_{\varphi} \left( f \right) \left( x \right) = f \left( x + \varphi \right)\)。证明它是线性变换，求其在基\(\sin x , \cos x\)下的矩阵、逆变换，并与平面旋转比较。
\solutionanswer
对任意\(f,g\in V\)与实数\(a,b\)，
\[
T_\varphi(af+bg)(x)=af(x+\varphi)+bg(x+\varphi),
\]
故满足线性性。由和角公式，
\[
\begin{aligned}
T_\varphi(\sin x)&=\cos\varphi\sin x+\sin\varphi\cos x,\\
T_\varphi(\cos x)&=-\sin\varphi\sin x+\cos\varphi\cos x.
\end{aligned}
\]
这也验证了像仍在\(V\)中。矩阵为
\[
A_\varphi=\begin{bmatrix}\cos\varphi&-\sin\varphi\\
\sin\varphi&\cos\varphi\end{bmatrix},
\]
逆变换是\(T_{-\varphi}\)。

在基\((\sin x,\cos x)\)给出的系数平面中，这个矩阵正是逆时针旋转\(\varphi\)的矩阵；在函数图像上，\(f(x+\varphi)\)则对应水平方向的平移。两种解释分别描述系数与自变量的变化。
\par
